% =====================================================================
%  Betriebsanweisung Feinschnitt-Tischkreissäge Proxxon FET
%  Eigenes PDF: Betriebsanweisung_Tischkreissaege.tex, Seite in der
%  Einweisung über \BAseite. Layout: fablab-document/fablab-ba.sty
%
%  Lizenz-Hinweis: Text vollständig selbst formuliert, KEINE Passagen und
%  KEINE Abbildungen aus der Proxxon-Betriebsanleitung übernehmen!
%  Sicherheitsregeln als solche (Fakten) sind frei, Wortlaut und Bilder
%  von Proxxon sind urheberrechtlich geschützt. Nur ISO-7010-Symbole.
% =====================================================================
\BAsetup{
  typ=maschine,
  titel={Feinschnitt-Tischkreissäge Proxxon FET},
  kurztitel={Tischkreissäge},
  taetigkeit={Längs-, Quer- und Gehrungsschnitte, Sägeblattwechsel, Reinigen},
  nummer=BA-TK-01,
  stand=TT.MM.JJJJ,                   % Datum der Freigabe
  verantwortlich={Name der verantwortlichen Person},
  signalwort=WARNUNG,
  schrift=\footnotesize,
  entwurf,                            % entfernen, sobald freigegeben
}

\BAkopf

\BAblock{ANWENDUNGSBEREICH}{M002}{%
  \begin{itemize}
    \item Feine, gerade Schnitte in dünnen Werkstücken (bis 22\,mm): Holz, Kunststoff, mit passendem Blatt auch Aluminium, Leiterplatten und Keramik.
    \item Benutzung nur nach unterschriebener Einweisung; Betriebsanleitung des Herstellers beachten.
  \end{itemize}}

\BAblock{GEFAHREN FÜR MENSCH UND UMWELT}{W022,W024}{%
  \begin{itemize}
    \item Schnittverletzungen an Fingern und Händen durch das laufende oder auslaufende Sägeblatt.
    \item \textbf{Rückschlag:} Klemmt das Werkstück, kann es zur Bedienperson geschleudert werden.
    \item Wegfliegende Späne und Splitter; Lärm.
    \item Gesundheitsschädlicher Staub (Hartholz, MDF, GFK-Leiterplatten).
    \item Einzugsgefahr für Handschuhe, weite Kleidung, Schmuck und lange Haare.
  \end{itemize}}

\BAblock{SCHUTZMASSNAHMEN UND VERHALTENSREGELN}{M003,M004,P028}{%
  \begin{itemize}
    \item Vor jeder Benutzung Säge, Kabel, Blatt, Spaltkeil und Sägeblattschutz prüfen; Defekte nicht benutzen und melden.
    \item Gehörschutz und Schutzbrille tragen. \textbf{Beim Sägen keine Handschuhe}, eng anliegende Kleidung, Haare zusammenbinden.
    \item Säge mit Zwingen auf der Werkbank befestigen, Absaugung anschließen (Sauger auf AUTO).
    \item Sägeblattschutz und Spaltkeil nie entfernen oder blockieren.
    \item Werkstück immer am Längs- oder Winkelanschlag führen, nie beide zugleich, nie freihändig.
    \item Blatt nur knapp über das Werkstück stellen. Hände nie in die Verlängerung des Blatts, kleine und schmale Teile mit dem Schiebestock schieben.
    \item Seitlich neben der Schnittlinie stehen. Reste erst entfernen, wenn das Blatt steht.
  \end{itemize}}

\BAblock{VERHALTEN BEI STÖRUNGEN UND IM GEFAHRENFALL}{W001,M006}{%
  \begin{itemize}
    \item Säge ausschalten, Blatt stehen lassen, Netzstecker ziehen, Betreuer informieren.
    \item Klemmendes Werkstück erst bei stehendem Blatt lösen.
    \item Entstehungsbrand: Säge vom Netz trennen, löschen. \textbf{Feuerwehr: 112}
  \end{itemize}}

\BAblock{VERHALTEN BEI UNFÄLLEN – ERSTE HILFE}{E003}{%
  \begin{itemize}
    \item Säge ausschalten, Stecker ziehen. Betreuer/Ersthelfer verständigen, Wunden versorgen.
    \item Jeden Unfall melden und im Verbandbuch eintragen.
  \end{itemize}\vspace{1.5mm}
  \textbf{Notruf: 112}\qquad FabLab: 09131\,/\,85\,28013}

\BAletzterblock{INSTANDHALTUNG, ENTSORGUNG}{M021}{%
  \begin{itemize}
    \item Sägeblattwechsel und Wartung nur durch Betreuer, mit gezogenem Netzstecker und Schutzhandschuhen.
    \item Nach der Arbeit Säge und Platz aussaugen, nicht mit Druckluft abblasen. Elektrische Prüfung nach DGUV Vorschrift 3 gemäß Prüffrist.
  \end{itemize}}
